\documentclass[10pt,a4paper]{amsart}
\usepackage[margin=19mm]{geometry}
\usepackage{amsmath,amssymb,mathtools}
\usepackage[scr=rsfs,cal=euler]{mathalfa}
\usepackage{xfrac}
\usepackage[unicode,hidelinks,bookmarks=false]{hyperref}
\newcommand{\ie}{i.e.\ }
\newcommand{\cf}{cf.\ }
\newcommand{\resp}{respectively\ }
\newcommand{\Subsection}{\subsection}
\providecommand{\nobreakdash}{\nobreak\hbox{-}}
\setlength{\emergencystretch}{2em}
\multlinegap0pt
\usepackage{tikz}
\usetikzlibrary{cd}
\usepackage{bm}
\usepackage{amssymb}
\usepackage{tikz}
\usepackage{xcolor}
\usepackage{enumerate}
\newtheorem{lemma}{Lemma}
\newtheorem{corollary}{Corollary}
\newcommand{\EE}{\mathcal{E}}
\newcommand{\PP}{\mathbb{P}}
\newcommand{\TP}{{\mathbb{T}\mathbb{P}}}
\newcommand{\TT}{\mathbb{T}}
\renewcommand{\phi}{\varphi}
\title{A tropical framework for using Porteous formula}
\author{Andrew R. Tawfeek}
\date{Source: arXiv:2411.10578v2 (CC BY 4.0)}
\begin{document}
\maketitle

\noindent\emph{Source and licence.} A tropical framework for using Porteous formula. Version arXiv:2411.10578v2 (CC BY 4.0), distributed by arXiv under the Creative Commons Attribution 4.0 International licence, http://creativecommons.org/licenses/by/4.0/, as recorded in the arXiv licence metadata for this exact version and shown on the abstract page https://arxiv.org/abs/2411.10578v2. Full licence notice: \texttt{licenses/arxiv-2411.10578-CC-BY-4.0.txt}.\par\smallskip

\noindent\textit{Editorial scope.} Counted unit: the article's corollary splitting (source0.tex lines 1017--1023). The article's complete original proof of it is delivered with it (source0.tex lines 1025--1058, one \texttt{proof} environment of 34 lines). No mathematics is written, repaired or re-derived: the statement and the proof are the article's own lines, in the article's own notation, and no retained line is substituted or re-flowed.  Retained uncounted as the source-local context the proof needs: lines 992--996.  Nothing was omitted from any retained span, so the package carries no omission record.  No retained line contains a citation, so the wrapper carries no bibliography and no bibliography entry.  The retained lines contain the article's own diagram code, which the wrapper typesets with the standard tikz packages; no image file is delivered and the package declares no asset dependency.  Editorial formatting only, in addition: an AMS \texttt{amsart} preamble that restates the article's own declarations and macros used by the retained lines, namely the environments \texttt{lemma}, \texttt{corollary} on separate counters, together with the standard packages those lines need.  The retained environments print in this document as Lemma 1, Corollary 1 in order of appearance, whereas the article prints the same items with the numbering of its own full document.  The complete native source of the article as distributed in the arXiv e-print for this exact version is bundled unchanged as \texttt{sources/168607/source0.tex}.
\begin{lemma} \label{splitpf}
    Let $\EE \to X$ be a rank $n+1$ tropical vector bundle. Then the induced pullback map
    $$\phi^*: A_k(X) \rightarrow A_{k+n}(\TT\PP(\EE))$$
    is a split monomorphism.
\end{lemma}

\begin{corollary}[Tropical Splitting Principle] \label{splitting}
For a tropical vector bundle $\pi: \mathcal{E} \to X$ of rank $n$ over a rational polyhedral space $X$, there is an associated rational polyhedral space $Y$ and a morphism $f: Y \to X$ such that
\begin{enumerate}
	\item the pullback bundle $f^*\mathcal{E}$ is a direct sum of tropical line bundles, and
	\item the induced pullback map $f^*: A(X) \to A(Y)$ is injective.
\end{enumerate}
\end{corollary}

\begin{proof}
We proceed with the classical \textit{splitting construction} adapted to our situation. At each step, $\TP(\EE_i)$ is a rational polyhedral space (with boundary from both the base and the fibers), and Lemma \ref{splitpf} applies via the projection formula. We construct $f$ by induction on the rank of $\EE$. Projectivization provides

\[\begin{tikzcd}
	{\varphi^*(\mathcal{E})} & {\mathcal{E}} \\
	{\mathbb{T}\mathbb{P}(\mathcal{E})} & X
	\arrow[""', dashed, from=1-1, to=2-1]
	\arrow["", dashed, from=1-1, to=1-2]
	\arrow["p", from=2-1, to=2-2]
	\arrow["\pi", from=1-2, to=2-2]
\end{tikzcd}\]

\noindent where the pullback $p^*\mathcal{E}$ fits into the short exact sequence
\begin{equation*}
0  \longrightarrow \mathcal{O}_{\TP (\mathcal{E})}(-1) \longrightarrow p^*\mathcal{E} \longrightarrow \mathcal{Q} \longrightarrow 0.
\end{equation*}
If $\mathcal{E}$ has rank $2$, then the quotient bundle $\mathcal{Q} = p^*(\mathcal{E})/\mathcal{O}_{\TP(\mathcal{E})} (-1)$ is a line bundle. Therefore we may take $Y = \TP(\EE)$ as $p^*\EE \cong \mathcal{O}_{\TP(\EE)} (-1) \oplus \mathcal{Q}$ and $p^*$ is injective by Lemma \ref{splitpf}. Otherwise, we may iterate this construction by repeated projectivization and quotienting:

% https://q.uiver.app/?q=WzAsOSxbNCwwLCJcXG1hdGhjYWx7RX1fMD1cXG1hdGhjYWx7RX0iXSxbNCwxLCJYIl0sWzMsMSwiXFx1bmRlcmJyYWNle1xcbWF0aGJie1R9XFxtYXRoYmJ7UH0oXFxtYXRoY2Fse0V9XzApfV97WV8xfSJdLFszLDAsIlxcb3ZlcmJyYWNle1xcdmFycGhpXzFeKihcXG1hdGhjYWx7RX1fMCkvXFxtYXRoY2Fse099X3t7XFxtYXRoY2Fse0V9fV8wfSgtMSl9XntcXG1hdGhjYWx7RX1fMX0iXSxbMiwxLCJcXHVuZGVyYnJhY2V7XFxtYXRoYmJ7VH1cXG1hdGhiYntQfShcXG1hdGhjYWx7RX1fMSl9X3tZXzJ9Il0sWzIsMCwiXFxvdmVyYnJhY2V7XFx2YXJwaGlfMl4qKFxcbWF0aGNhbHtFfV8xKS9cXG1hdGhjYWx7T31fe3tcXG1hdGhjYWx7RX19XzF9KC0xKX1ee1xcbWF0aGNhbHtFfV8yfSJdLFsxLDEsIlxcY2RvdHMiXSxbMCwxLCJcXHVuZGVyYnJhY2V7XFxtYXRoYmJ7VH1cXG1hdGhiYntQfShcXG1hdGhjYWx7RX1fbil9X3tZX259Il0sWzAsMCwiXFxvdmVyYnJhY2V7XFx2YXJwaGlfbl4qKFxcbWF0aGNhbHtFfV97bi0xfSkvXFxtYXRoY2Fse099X3t7XFxtYXRoY2Fse0V9fV97bi0xfX0oLTEpfV57XFxtYXRoY2Fse0V9X259Il0sWzAsMV0sWzIsMSwiXFx2YXJwaGlfMSIsMl0sWzMsMl0sWzQsMiwiXFx2YXJwaGlfMiIsMl0sWzYsNCwiXFx2YXJwaGlfMyIsMl0sWzUsNF0sWzcsNiwiXFx2YXJwaGlfbiIsMl0sWzgsN11d
\[\begin{tikzcd}[column sep=small]
	{\overbrace{\varphi_n^*(\mathcal{E}_{n-1})/\mathcal{O}_{{\mathcal{E}}_{n-1}}(-1)}^{\mathcal{E}_n}} && {\overbrace{\varphi_2^*(\mathcal{E}_1)/\mathcal{O}_{{\mathcal{E}}_1}(-1)}^{\mathcal{E}_2}} & {\overbrace{\varphi_1^*(\mathcal{E}_0)/\mathcal{O}_{{\mathcal{E}}_0}(-1)}^{\mathcal{E}_1}} & {\mathcal{E}_0=\mathcal{E}} \\
	{\underbrace{\mathbb{T}\mathbb{P}(\mathcal{E}_n)}_{Y_n}} & \cdots & {\underbrace{\mathbb{T}\mathbb{P}(\mathcal{E}_1)}_{Y_2}} & {\underbrace{\mathbb{T}\mathbb{P}(\mathcal{E}_0)}_{Y_1}} & X
	\arrow[from=1-5, to=2-5]
	\arrow["{p_1}"', from=2-4, to=2-5]
	\arrow[from=1-4, to=2-4]
	\arrow["{p_2}"', from=2-3, to=2-4]
	\arrow["{p_3}"', from=2-2, to=2-3]
	\arrow[from=1-3, to=2-3]
	\arrow["{p_n}"', from=2-1, to=2-2]
	\arrow[from=1-1, to=2-1]
\end{tikzcd}\]
Then taking $Y = Y_n$, we have that the pullback decomposes into a direct sum  $f^*\EE = \bigoplus_{j=1}^n \mathcal{L}_j$, and that $f^*: A(X) \to A(Y)$ is the desired injective map. 

\end{proof}

\end{document}
